\documentclass[12pt,a4paper]{article}

% ── Packages ──
\usepackage[margin=1in]{geometry}
\usepackage{amsmath,amssymb,amsfonts}
\usepackage{booktabs}
\usepackage{graphicx}
\usepackage{hyperref}
\usepackage{natbib}
\usepackage{float}
\usepackage{enumitem}
\usepackage{caption}
\usepackage{subcaption}
\usepackage{xcolor}
\usepackage{multirow}
\usepackage{array}
\usepackage{longtable}
\usepackage{setspace}

\hypersetup{
    colorlinks=true,
    linkcolor=blue!60!black,
    citecolor=blue!60!black,
    urlcolor=blue!60!black
}

\onehalfspacing

\title{%
    \textbf{Cross-Sectional Factor Investing with Portfolio Optimization and Black-Litterman Allocation} \\[0.5em]
    \large An End-to-End Quantitative Investment Pipeline
}
\author{Zian (Andy) Wang}
\date{March 2026}

\begin{document}

\maketitle

\begin{abstract}
We build an end-to-end quantitative factor investment pipeline that constructs 20 cross-sectional equity factors from Compustat/CRSP data (1970--2019), selects a parsimonious subset via four complementary methods, and allocates across factor portfolios using mean-variance optimization, risk parity, Black-Litterman posterior returns, and MAXSER sparse regression for joint factor+stock portfolios. Out-of-sample (2015--2019), all nine portfolio variants achieve Sharpe ratios above 0.93, with Equal Weight reaching 1.20 on a market-neutral basis, compared to 0.91 for the S\&P~500 and 0.56 for a hedge fund index composite. MAXSER factor+stock variants achieve Sharpe ratios of 1.13 with 14.7\% annualized returns. Rolling backtests with 10~bps transaction costs confirm robustness: Equal Weight maintains a net Sharpe of 1.21 with only 5\% turnover. Black-Litterman improves MVO's net Sharpe from 0.76 to 1.00 while halving turnover, while adaptive annual factor re-selection underperforms fixed selection across all strategies---demonstrating the cost of factor chasing.
\end{abstract}

\tableofcontents
\newpage

% ═══════════════════════════════════════════════════════════════
\section{Introduction}
% ═══════════════════════════════════════════════════════════════

Factor investing is the systematic construction of portfolios that harvest risk premia associated with firm characteristics such as value, momentum, quality, and size. Since the seminal work of \citet{fama1993common} and \citet{carhart1997persistence}, academic research has identified dozens of cross-sectional predictors of stock returns, though many fail to survive out-of-sample testing or transaction costs \citep{harvey2016and,hou2020replicating}.

This project addresses three questions:
\begin{enumerate}[noitemsep]
    \item Which factors from a broad set of 20 candidates deliver robust, uncorrelated alpha?
    \item How should capital be allocated across selected factors---and does Black-Litterman improve upon na\"ive optimization?
    \item Do the resulting portfolios survive rolling-window backtests with realistic transaction costs?
\end{enumerate}

We build a fully automated, reproducible pipeline that takes raw Compustat/CRSP data as input and produces optimized, backtested factor portfolios as output. The pipeline is implemented in Python and organized into modular stages, each of which caches intermediate results for independent re-execution.

% ═══════════════════════════════════════════════════════════════
\section{Data}
% ═══════════════════════════════════════════════════════════════

\subsection{Universe}

Our primary dataset is the CRSP/Compustat merged file covering all U.S.\ common stocks (share codes 10, 11) from January 1970 to December 2019. We filter to S\&P~500 constituents at each month-end to ensure investability and liquidity.

\subsection{Sample Periods}

\begin{itemize}[noitemsep]
    \item \textbf{Full sample}: January 1970 -- December 2019 (factor construction)
    \item \textbf{Validation}: January 1987 -- December 2019 (quintile sort analysis)
    \item \textbf{In-sample (IS)}: January 1987 -- December 2014 (optimization, selection)
    \item \textbf{Out-of-sample (OOS)}: January 2015 -- December 2019 (evaluation)
\end{itemize}

\subsection{Benchmarks}

We compare against 14 benchmarks spanning passive indices, Fama-French factors, mutual funds, hedge fund indices, and smart beta ETFs:

\begin{itemize}[noitemsep]
    \item \textbf{Market}: S\&P~500 total return, Fama-French Mkt-RF
    \item \textbf{Fama-French}: SMB, HML, UMD
    \item \textbf{Mutual funds}: FSMEX, FLPSX, FCNTX, FSPHX, FBALX (equal-weight and best)
    \item \textbf{Hedge fund indices}: HFRI composite strategies (equal-weight and best)
    \item \textbf{Smart beta}: Factor-tilted ETFs (equal-weight and best)
\end{itemize}

All benchmark Sharpe ratios are computed on excess returns (subtracting the risk-free rate). Our long-short factor portfolios are market-neutral by construction, so no risk-free adjustment is needed.

% ═══════════════════════════════════════════════════════════════
\section{Factor Construction}
\label{sec:factors}
% ═══════════════════════════════════════════════════════════════

We construct 20 cross-sectional factors spanning seven categories (Table~\ref{tab:factors}). Each factor is computed monthly from lagged accounting data (to avoid look-ahead bias) and winsorized at the 1st and 99th percentiles.

\begin{table}[H]
\centering
\caption{Factor Universe}
\label{tab:factors}
\small
\begin{tabular}{lll}
\toprule
\textbf{Category} & \textbf{Factor} & \textbf{Definition} \\
\midrule
\multirow{4}{*}{Value}
    & Book-to-Price (BP) & Book equity / market cap \\
    & Cash Flow to Price (CFTP) & Operating cash flow / market cap \\
    & Earnings Yield & Net income / market cap \\
    & Dividend Yield & Dividends per share / price \\
\midrule
Size & Log Market Cap & $\ln(\text{price} \times \text{shares outstanding})$ \\
\midrule
\multirow{4}{*}{Momentum}
    & 12-1 Momentum (MOM) & Cumulative return months $t{-}12$ to $t{-}2$ \\
    & Short-Term Reversal (STReversal) & Prior month return \\
    & High-Low 1M (HL1M) & Monthly high/low price ratio \\
    & Long-Term Growth (LTGC) & Long-term price reversal \\
\midrule
\multirow{3}{*}{Quality}
    & Accrual Ratio & $(\Delta\text{CA} - \Delta\text{Cash} - \Delta\text{CL}) / \text{TA}$ \\
    & Return on Equity (ROE) & Net income / book equity \\
    & Net Profit Margin & Net income / revenue \\
\midrule
\multirow{3}{*}{Growth}
    & Sales Growth & YoY revenue growth \\
    & Asset Growth & YoY total asset growth \\
    & Sustainable Growth & $\text{ROE} \times (1 - \text{payout ratio})$ \\
\midrule
\multirow{2}{*}{Analyst}
    & Earnings Revision & Consensus estimate changes \\
    & SUE & Standardized unexpected earnings \\
\midrule
Risk & Ann.\ Volatility (12M) & Trailing 12-month return std $\times \sqrt{12}$ \\
\midrule
Leverage & Debt-to-Equity & Total debt / book equity \\
\midrule
Other & Gross Profit & Gross profit / total assets \\
\bottomrule
\end{tabular}
\end{table}

\subsection{Quintile Sort Methodology}

For each factor and each month, we sort S\&P~500 stocks into quintiles (Q1 = lowest, Q5 = highest). The \textbf{QSpread} is the equal-weighted return differential between the long and short quintiles, with the sign chosen so that a positive QSpread indicates the factor delivers positive alpha (e.g., for value factors, long Q5/short Q1; for growth factors, long Q1/short Q5).

\subsection{Benchmark Validation}

We validate our factor construction against Capital IQ factor indices. Correlations between our QSpreads and Capital IQ benchmark returns range from \textbf{0.87 to 0.99} across 8 matched factors:

\begin{table}[H]
\centering
\caption{Factor Validation: Correlation with Capital IQ Benchmarks}
\label{tab:validation}
\begin{tabular}{lc}
\toprule
\textbf{Factor} & \textbf{Correlation} \\
\midrule
Log Market Cap & 0.993 \\
HL1M & 0.992 \\
Book-to-Price & 0.989 \\
Annualized Volatility & 0.965 \\
Momentum (12-1) & 0.925 \\
Long-Term Growth & 0.915 \\
Net Profit Margin & 0.894 \\
ROE & 0.876 \\
\bottomrule
\end{tabular}
\end{table}

% ═══════════════════════════════════════════════════════════════
\section{Factor Selection}
\label{sec:selection}
% ═══════════════════════════════════════════════════════════════

From the 20-factor universe, we select a parsimonious subset that balances signal strength with diversification. We employ four complementary selection methods and combine their votes:

\subsection{Fama-MacBeth Regression}

For each month $t$ in the in-sample period, we run cross-sectional regressions of stock returns on factor exposures:
\begin{equation}
    r_{i,t} = \alpha_t + \sum_{k} \gamma_{k,t} \cdot f_{i,k,t-1} + \epsilon_{i,t}
\end{equation}
The time-series average $\bar{\gamma}_k$ is tested for significance using Newey-West standard errors (6 lags). Factors with $|t| > 2$ are considered significant.

\subsection{Information Coefficient Analysis}

The Information Coefficient (IC) is the rank correlation between factor values and subsequent returns:
\begin{equation}
    \text{IC}_{k,t} = \text{Spearman}\!\left(f_{i,k,t-1},\; r_{i,t}\right)
\end{equation}
We filter on $|\overline{\text{IC}}| > 0.02$ and $\text{IC-IR} = \overline{\text{IC}} / \sigma(\text{IC}) > 0.5$.

\subsection{LASSO Selection}

We fit a LASSO regression with 5-fold time-series cross-validation:
\begin{equation}
    \hat{\boldsymbol{\gamma}} = \arg\min_{\boldsymbol{\gamma}} \left\| \mathbf{r} - \mathbf{F}\boldsymbol{\gamma} \right\|_2^2 + \lambda \left\| \boldsymbol{\gamma} \right\|_1
\end{equation}
Factors with nonzero coefficients at the optimal $\lambda$ are selected.

\subsection{Greedy Forward Selection}

Starting from an empty set, we iteratively add the factor with the highest $|\text{QSpread Sharpe}|$, subject to the constraint that all pairwise correlations with existing factors remain below 0.6:
\begin{equation}
    f^* = \arg\max_{f \notin S} \left| \text{Sharpe}(Q_f) \right| \quad \text{s.t.} \quad |\rho(Q_f, Q_s)| < 0.6 \;\; \forall s \in S
\end{equation}

\subsection{Selected Factors}

Consensus across the four methods yields five factors:

\begin{table}[H]
\centering
\caption{Selected Factor Combination}
\label{tab:selected}
\begin{tabular}{llcc}
\toprule
\textbf{Factor} & \textbf{Category} & \textbf{IS QSpread Sharpe} & \textbf{Max Pairwise $|\rho|$} \\
\midrule
Accrual Ratio & Quality & 0.84 & --- \\
CFTP & Value & 0.72 & 0.55 \\
STReversal & Momentum & 0.68 & 0.21 \\
Asset Growth & Growth & 0.61 & 0.38 \\
ROE & Quality & 0.58 & 0.45 \\
\bottomrule
\end{tabular}
\end{table}

The maximum pairwise correlation among selected factors is 0.55 (AccrualRatio--CFTP), ensuring genuine diversification across alpha sources.

% ═══════════════════════════════════════════════════════════════
\section{Portfolio Optimization}
\label{sec:optimization}
% ═══════════════════════════════════════════════════════════════

We allocate capital across the five selected factor QSpreads using multiple optimization approaches. Let $\boldsymbol{\mu}$ and $\boldsymbol{\Sigma}$ denote the in-sample mean return vector and covariance matrix of the factor QSpreads.

\subsection{Equal Weight}

$w_k = 1/K$ for all $k = 1, \ldots, K$. The 1/N portfolio serves as a robust baseline that is difficult to beat out of sample \citep{demiguel2009optimal}.

\subsection{IC-Weighted}

Weights proportional to in-sample information coefficients:
\begin{equation}
    w_k = \frac{|\overline{\text{IC}}_k|}{\sum_{j} |\overline{\text{IC}}_j|}
\end{equation}

\subsection{Mean-Variance Optimization (MVO)}

\begin{equation}
    \mathbf{w}^* = \arg\min_{\mathbf{w}} \left( \mathbf{w}'\boldsymbol{\Sigma}\mathbf{w} - \frac{1}{\delta}\boldsymbol{\mu}'\mathbf{w} \right) \quad \text{s.t.} \quad \sum |w_k| \leq L_{\text{gross}}, \; w_k \in [w_{\min}, w_{\max}]
\end{equation}
with risk aversion $\delta = 10$, $w_{\min} = -2$, $w_{\max} = 2$, and gross leverage $L_{\text{gross}} = 3$.

\subsection{Maximum Sharpe Ratio}

\begin{equation}
    \mathbf{w}^* = \arg\max_{\mathbf{w}} \frac{\boldsymbol{\mu}'\mathbf{w}}{\sqrt{\mathbf{w}'\boldsymbol{\Sigma}\mathbf{w}}} \quad \text{s.t.} \quad \sum |w_k| \leq 3
\end{equation}

\subsection{Risk Parity}

Equal risk contribution from each factor:
\begin{equation}
    w_k \cdot (\boldsymbol{\Sigma}\mathbf{w})_k = \frac{1}{K} \cdot \mathbf{w}'\boldsymbol{\Sigma}\mathbf{w} \quad \forall k
\end{equation}

\subsection{Black-Litterman Variants}

We use Black-Litterman (BL) as a \textbf{return estimation method}: compute posterior expected returns, then feed to MVO or Max Sharpe. The BL posterior is:
\begin{equation}
    \bar{\boldsymbol{\mu}} = \left[(\tau\boldsymbol{\Sigma})^{-1} + \mathbf{P}'\boldsymbol{\Omega}^{-1}\mathbf{P}\right]^{-1} \left[(\tau\boldsymbol{\Sigma})^{-1}\boldsymbol{\pi} + \mathbf{P}'\boldsymbol{\Omega}^{-1}\mathbf{Q}\right]
\end{equation}
where $\boldsymbol{\pi} = \delta\boldsymbol{\Sigma}\mathbf{w}_{\text{eq}}$ is the equilibrium return (equal-weight prior), $\mathbf{P} = \mathbf{I}$ (absolute views on each factor), $\mathbf{Q} = \boldsymbol{\mu}_{\text{IS}}$ (in-sample mean returns as views), and $\boldsymbol{\Omega}$ is diagonal with $\Omega_{kk} = \tau \cdot \Sigma_{kk}$ following \citet{he1999intuition}.

This yields two additional variants: \textbf{MVO + BL} and \textbf{Max Sharpe + BL}.

\subsection{Factor+Stock Portfolios (MAXSER)}

In addition to factor-only allocation, we construct \textbf{factor+stock portfolios} that combine factor QSpread weights with stock-level idiosyncratic alpha using the MAXSER framework of \citet{ao2019approaching}. The key motivation is the \textbf{weight double-counting problem}: stock returns contain factor exposures ($r_i = \beta' f + \alpha_i + \varepsilon_i$), so applying standard optimizers (MVO, Max Sharpe) to the joint $[K \text{ factors} + N \text{ stocks}]$ universe double-counts factor risk and produces uninterpretable weights with inflated effective leverage. MAXSER solves this by decomposing the portfolio into two legs: (1)~a \textbf{factor leg} (tangency portfolio over $K$ factors, where sample estimates are reliable since $K$ is small), and (2)~an \textbf{idiosyncratic leg} (sparse regression on stock residuals after projecting out factor exposures), with beta-adjustment to ensure factor exposure is not counted twice. The sparse regression estimates the maximum squared Sharpe ratio $\hat{\theta}^2 = \boldsymbol{\mu}'\boldsymbol{\Sigma}^{-1}\boldsymbol{\mu}$ via Lasso and Ridge regularization. Stock subpool selection uses idiosyncratic information ratio (residual alpha per unit residual risk) to select 50 stocks where $N \ll T$ is satisfied.

\subsection{Portfolio Variant Space}

The full combinatorial space follows a four-level decision tree:
\begin{enumerate}[noitemsep]
    \item \textbf{Asset universe}: factor-only ($K=5$) or factor+stocks ($K + N$)
    \item \textbf{Optimizer}: Equal Weight, IC-Weighted, MVO, Max Sharpe, or Risk Parity
    \item \textbf{Return estimation}: raw sample $\boldsymbol{\mu}$ or BL posterior $\bar{\boldsymbol{\mu}}$
    \item \textbf{Estimation method} (factor+stocks only): MAXSER (Lasso/Ridge)
\end{enumerate}

This yields $5 \times 2 = 10$ factor-only and $5 \times 2 \times 2 = 20$ factor+stock combinations, for 30 theoretical variants. We implement 9 and exclude 21 for principled reasons (Table~\ref{tab:exclusions}).

\begin{table}[H]
\centering
\caption{Excluded Portfolio Combinations and Rationale}
\label{tab:exclusions}
\small
\begin{tabular}{p{4.5cm}p{10cm}}
\toprule
\textbf{Excluded Combo} & \textbf{Reason} \\
\midrule
EW + BL (factor-only) & Equal weight assigns $1/K$ regardless of returns---BL posterior $\bar{\boldsymbol{\mu}}$ has no effect on the allocation. \\
IC-Weighted + BL (all) & IC weights derive from rank correlation magnitude, not expected returns---BL adjusts $\boldsymbol{\mu}$ but IC-Weighted does not use $\boldsymbol{\mu}$. \\
Risk Parity + BL (all) & Risk Parity uses only $\boldsymbol{\Sigma}$, never $\boldsymbol{\mu}$. While BL does produce a posterior $\boldsymbol{\Sigma}$, the shift is small ($\tau\boldsymbol{\Sigma}$ adjustment) and Risk Parity's purpose is to avoid return estimation entirely. \\
EW (factor+stocks) & Assigning $1/(K+N)$ across ${\sim}279$ assets gives each stock ${\sim}0.4\%$ weight, drowning factor signal in stock-level noise. \\
Risk Parity (factor+stocks) & Same issue: hundreds of stocks dominate the risk budget over 5 factors. \\
Plug-in F+S (MVO, MaxSharpe, IC-Weighted) & Factors are portfolios of stocks---optimizing over both in a joint $[K+N]$ universe double-counts exposure and inflates effective leverage. MAXSER avoids this via its own sparse regression framework. \\
MAXSER + BL & MAXSER includes its own shrinkage via Lasso/Ridge regularization on the squared Sharpe estimator. Layering BL posterior returns would double-regularize and conflate two estimation frameworks. \\
MAXSER + specific optimizer & MAXSER \textit{is} its own optimizer (maximizes $\hat{\theta}^2 = \boldsymbol{\mu}'\boldsymbol{\Sigma}^{-1}\boldsymbol{\mu}$ via sparse regression)---it replaces MVO/Max Sharpe rather than composing with them. \\
\bottomrule
\end{tabular}
\end{table}

% ═══════════════════════════════════════════════════════════════
\section{Out-of-Sample Results}
\label{sec:results}
% ═══════════════════════════════════════════════════════════════

Table~\ref{tab:oos} reports OOS performance (January 2015 -- December 2019) for all portfolio variants and benchmarks. Our factor-only portfolios are market-neutral long-short (long Q1, short Q5 within S\&P~500), so absolute returns reflect pure cross-sectional alpha with zero market beta. MAXSER factor+stock variants combine factor QSpreads with stock-level idiosyncratic alpha via sparse regression, producing higher absolute returns with commensurately higher volatility.

\begin{table}[H]
\centering
\caption{Out-of-Sample Performance Comparison (2015--2019)}
\label{tab:oos}
\footnotesize
\resizebox{\textwidth}{!}{%
\begin{tabular}{lccccccc}
\toprule
\textbf{Portfolio} & \textbf{Sharpe} & \textbf{Sortino} & \textbf{Ann.\ Ret.} & \textbf{Ann.\ Vol} & \textbf{Max DD} & \textbf{VaR\textsubscript{95}} & \textbf{CF-VaR\textsubscript{95}} \\
\midrule
\multicolumn{8}{l}{\textit{Our Portfolios --- Factor-Only (market-neutral long-short)}} \\
Equal Weight    & \textbf{1.203} & 2.464 & 3.9\%  & 3.2\%  & $-$3.3\%  & $-$1.1\% & $-$1.2\% \\
IC-Weighted     & 1.122          & 2.298 & 3.9\%  & 3.5\%  & $-$2.2\%  & $-$1.1\% & $-$1.2\% \\
MVO + BL        & 1.114          & 2.062 & 5.5\%  & 4.9\%  & $-$5.1\%  & $-$1.6\% & $-$1.7\% \\
MVO             & 1.038          & 1.861 & 8.0\%  & 7.7\%  & $-$8.3\%  & $-$2.9\% & $-$2.7\% \\
Max Sharpe + BL & 1.023          & 1.778 & 7.8\%  & 7.7\%  & $-$9.4\%  & $-$2.8\% & $-$2.6\% \\
Risk Parity     & 0.998          & 1.814 & 3.1\%  & 3.1\%  & $-$2.8\%  & $-$1.2\% & $-$1.2\% \\
Max Sharpe      & 0.938          & 1.661 & 10.0\% & 10.6\% & $-$12.8\% & $-$4.3\% & $-$3.7\% \\
\midrule
\multicolumn{8}{l}{\textit{Our Portfolios --- Factor+Stock (MAXSER)}} \\
MAXSER-Lasso    & 1.129          & 2.049 & 14.7\% & 13.0\% & $-$13.8\% & $-$4.8\% & $-$4.7\% \\
MAXSER-Ridge    & 1.125          & 2.213 & 14.6\% & 13.0\% & $-$13.2\% & $-$4.8\% & $-$4.3\% \\
\midrule
\multicolumn{8}{l}{\textit{Benchmarks}} \\
HFRIMAI (Best HF)  & 1.358 & 2.371 & 3.2\%  & 2.4\%  & $-$2.4\%  & $-$1.0\% & $-$1.0\% \\
Mutual Fund (EW)   & 0.947 & 1.357 & 16.4\% & 17.4\% & $-$21.1\% & $-$6.3\% & $-$6.7\% \\
Mutual Fund (FSMEX)& 0.919 & 1.465 & 19.0\% & 20.7\% & $-$17.1\% & $-$8.8\% & $-$8.2\% \\
S\&P 500           & 0.907 & 1.219 & 10.8\% & 11.9\% & $-$14.1\% & $-$6.0\% & $-$5.2\% \\
FF Mkt-RF          & 0.857 & 1.103 & 10.7\% & 12.4\% & $-$15.1\% & $-$6.1\% & $-$5.5\% \\
Smart Beta (PSL)   & 0.818 & 1.059 & 8.0\%  & 9.8\%  & $-$10.6\% & $-$4.0\% & $-$4.6\% \\
Smart Beta (EW)    & 0.744 & 0.925 & 8.9\%  & 11.9\% & $-$16.4\% & $-$5.4\% & $-$5.3\% \\
HF Index (EW)      & 0.562 & 0.707 & 2.0\%  & 3.6\%  & $-$7.8\%  & $-$1.8\% & $-$1.7\% \\
FF Momentum (UMD)  & 0.201 & 0.351 & 2.6\%  & 13.2\% & $-$21.3\% & $-$6.3\% & $-$5.9\% \\
FF SMB             & $-$0.257 & $-$0.547 & $-$2.2\% & 8.6\% & $-$17.2\% & $-$3.6\% & $-$4.0\% \\
FF HML             & $-$0.472 & $-$0.968 & $-$4.3\% & 9.0\% & $-$30.1\% & $-$3.9\% & $-$3.8\% \\
\bottomrule
\end{tabular}}
\end{table}

\textbf{Key observations:}
\begin{itemize}[noitemsep]
    \item All nine variants achieve Sharpe $\geq 0.94$, with eight exceeding the S\&P~500 (0.91) on a risk-adjusted basis. Only HFRIMAI---a low-volatility hedge fund sub-strategy index---beats our best variant.
    \item \textbf{Market neutrality explains low absolute returns.} Factor-only portfolios earn 3--10\% annually from pure cross-sectional alpha with zero beta. The S\&P~500's 10.8\% is mostly equity risk premium. Our max drawdowns (2--13\%) are 2--5$\times$ smaller than benchmarks (14--21\%).
    \item \textbf{Simple allocations are hard to beat}: Equal Weight (Sharpe 1.20) outperforms all optimized variants, echoing \citet{demiguel2009optimal}. IC-Weighted (1.12) and MVO+BL (1.11) follow.
    \item \textbf{MAXSER captures idiosyncratic alpha}: MAXSER-Lasso (Sharpe 1.13, return 14.7\%) and MAXSER-Ridge (1.13, 14.6\%) achieve the highest absolute returns by combining factor allocation with stock-level residual alpha from 50 selected stocks.
    \item MVO + BL (Sharpe 1.11) outperforms plain MVO (1.04), showing that BL regularization helps by shrinking returns toward equilibrium.
    \item The Fama-French value factor (HML) has a \textit{negative} Sharpe ($-$0.47) during 2015--2019, highlighting the ``value winter'' that plagued traditional value strategies.
\end{itemize}

\subsection{Cornish-Fisher VaR}

We report both historical VaR and Cornish-Fisher (CF) adjusted VaR at the 95\% level. The CF adjustment accounts for non-normal skewness and kurtosis:
\begin{equation}
    z_{\text{CF}} = z_\alpha + \frac{1}{6}(z_\alpha^2 - 1)S + \frac{1}{24}(z_\alpha^3 - 3z_\alpha)(K-3) - \frac{1}{36}(2z_\alpha^3 - 5z_\alpha)S^2
\end{equation}
where $S$ is skewness, $K$ is kurtosis, and $z_\alpha$ is the standard normal quantile. CF-VaR generally stays close to historical VaR for our factor portfolios (returns are approximately symmetric), but shows larger deviations for benchmarks with fat tails (e.g., mutual funds, S\&P~500).

\subsection{Sharpe Ratio Equality Tests}

We test whether our portfolios statistically outperform benchmarks using the Jobson-Korkie test with Memmel correction. For Equal Weight vs S\&P~500, the difference in Sharpe ratios (1.20 vs 0.91) is economically meaningful, though statistical significance depends on the short OOS window (60 months). The test has low power with $T=60$ monthly observations---a longer OOS period or bootstrap confidence intervals would strengthen inference.

% ═══════════════════════════════════════════════════════════════
\section{Black-Litterman Robustness Analysis}
\label{sec:bl}
% ═══════════════════════════════════════════════════════════════

Stage~3 uses BL as a return estimator inside MVO and Max Sharpe (Section~\ref{sec:optimization}). This section tests whether that improvement is robust to the choice of hyperparameters ($\tau$, $\delta$) and whether it is statistically significant.

\subsection{Hyperparameter Sensitivity}

The baseline $\tau = 0.05$, $\delta = 10$ was chosen \textit{a priori} from \citet{he1999intuition}. We evaluate OOS Sharpe ratios across a grid of $\tau \in \{0.01, 0.025, 0.05, 0.1, 0.25, 0.5\}$ and $\delta \in \{2.5, 5, 10, 25, 50\}$, re-running the BL posterior through the \textit{exact same} Stage~3 optimizers (MVO with leverage $\leq 2\times$, Max Sharpe with leverage $\leq 3\times$).

\textbf{Important}: This grid uses OOS data for \textit{evaluation}, not parameter selection. Choosing $\tau$/$\delta$ to maximize OOS Sharpe would be data snooping.

\textbf{Tau cancellation.} Under He \& Litterman's recommended $\boldsymbol{\Omega}$ scaling ($\Omega_{kk} = \tau \cdot \Sigma_{kk} / c_k$), $\tau$ appears in both the prior precision $(\tau\boldsymbol{\Sigma})^{-1}$ and the view precision $\boldsymbol{\Omega}^{-1}$, canceling out of the posterior mean $\bar{\boldsymbol{\mu}}$. Since the optimizers use only $\bar{\boldsymbol{\mu}}$ (not the posterior covariance), $\tau$ has no effect. We fix $\boldsymbol{\Omega}$ at the baseline $\tau = 0.05$ to break this cancellation and make $\tau$ sensitivity meaningful.

\begin{table}[H]
\centering
\caption{MVO + BL OOS Sharpe Across ($\tau$, $\delta$) Grid}
\label{tab:bl_grid_mvo}
\small
\begin{tabular}{lccccc}
\toprule
 & $\delta=2.5$ & $\delta=5$ & $\delta=10$ & $\delta=25$ & $\delta=50$ \\
\midrule
$\tau=0.01$  & 0.76 & 0.97 & \textbf{1.17} & \textbf{1.20} & \textbf{1.16} \\
$\tau=0.025$ & 0.75 & 0.77 & \textbf{1.02} & \textbf{1.13} & \textbf{1.08} \\
$\tau=0.05$  & 0.75 & 0.74 & 0.96 & \textbf{1.08} & \textbf{1.03} \\
$\tau=0.1$   & 0.76 & 0.74 & 0.95 & \textbf{1.06} & \textbf{1.00} \\
$\tau=0.25$  & 0.76 & 0.74 & 0.95 & \textbf{1.05} & 0.98 \\
$\tau=0.5$   & 0.77 & 0.74 & 0.95 & \textbf{1.06} & 0.98 \\
\bottomrule
\end{tabular}
\\[0.3em]
\small Plain MVO baseline: 0.97. Bold = BL beats plain MVO. BL improves MVO in 50\% of grid cells.
\end{table}

\begin{table}[H]
\centering
\caption{Max Sharpe + BL OOS Sharpe Across ($\tau$, $\delta$) Grid}
\label{tab:bl_grid_ms}
\small
\begin{tabular}{lccccc}
\toprule
 & $\delta=2.5$ & $\delta=5$ & $\delta=10$ & $\delta=25$ & $\delta=50$ \\
\midrule
$\tau=0.01$  & 0.75 & \textbf{0.94} & \textbf{1.17} & \textbf{1.30} & \textbf{1.27} \\
$\tau=0.025$ & 0.73 & 0.79 & \textbf{0.94} & \textbf{1.21} & \textbf{1.27} \\
$\tau=0.05$  & 0.75 & 0.78 & 0.84 & \textbf{1.07} & \textbf{1.23} \\
$\tau=0.1$   & 0.79 & 0.80 & 0.84 & \textbf{0.93} & \textbf{1.11} \\
$\tau=0.25$  & 0.83 & 0.83 & 0.85 & \textbf{0.88} & \textbf{0.94} \\
$\tau=0.5$   & 0.85 & 0.85 & 0.86 & \textbf{0.87} & \textbf{0.90} \\
\bottomrule
\end{tabular}
\\[0.3em]
\small Plain Max Sharpe baseline: 0.87. Bold = BL beats plain Max Sharpe. BL improves Max Sharpe in 50\% of grid cells.
\end{table}

The pattern is clear: BL helps most at \textbf{higher $\delta$} (more risk aversion shrinks weights toward equilibrium) and \textbf{lower $\tau$} (heavier prior weight stabilizes return estimates). The best region ($\tau \leq 0.05$, $\delta \geq 25$) consistently delivers Sharpe $> 1.0$. Our baseline ($\tau = 0.05$, $\delta = 10$) is conservative---it captures moderate BL improvement without overfitting to the grid.

\subsection{Statistical Significance}

We test whether the BL improvement is statistically significant using the Jobson-Korkie test (Sharpe ratio equality with Memmel correction) and the Diebold-Mariano test (equal predictive accuracy):

\begin{table}[H]
\centering
\caption{Statistical Tests of BL Improvement (Baseline $\tau=0.05$, $\delta=10$)}
\label{tab:bl_tests}
\small
\begin{tabular}{lcccc}
\toprule
\textbf{Comparison} & \textbf{$\Delta$ Sharpe} & \textbf{JK $z$} & \textbf{JK $p$} & \textbf{Significant?} \\
\midrule
MVO+BL vs MVO             & +0.07 & 0.62 & 0.53 & No \\
MaxSharpe+BL vs MaxSharpe & +0.09 & 1.41 & 0.16 & No \\
\bottomrule
\end{tabular}
\end{table}

Neither test rejects $H_0$ at the 5\% level. With only $T = 60$ monthly OOS observations, the Jobson-Korkie test has low power for detecting Sharpe differences of $\sim$0.1. The \textit{economic} significance may be more informative: BL consistently improves Sharpe and halves turnover (Section~\ref{sec:backtest}), and the sensitivity grid shows improvement across a wide region of the parameter space.

\subsection{Stock-Level BL}

For completeness, we implement stock-level BL on 274 S\&P~500 stocks. The stock-level BL (Sharpe 0.32) underperforms market-cap weighting (0.83), as views based on cross-sectional factor scores are too diffuse when applied to individual stocks. This motivates our focus on factor-level allocation: with only $K=5$ factors, each view carries meaningful information, whereas spreading views across $N \approx 274$ stocks dilutes signal into noise.

% ═══════════════════════════════════════════════════════════════
\section{Rolling Backtest with Transaction Costs}
\label{sec:backtest}
% ═══════════════════════════════════════════════════════════════

Static optimization on a single IS/OOS split can overstate performance. We run a rolling monthly backtest to evaluate robustness:

\begin{itemize}[noitemsep]
    \item \textbf{Lookback}: 36 months (expanding from the start of the validation period)
    \item \textbf{Rebalance frequency}: Quarterly (every 3 months)
    \item \textbf{Transaction costs}: 10 bps per unit turnover
    \item \textbf{Covariance estimation}: Ledoit-Wolf shrinkage at each rebalance
    \item \textbf{Weight drift}: Weights drift with returns between rebalance dates
\end{itemize}

\subsection{Results}

\begin{table}[H]
\centering
\caption{Rolling Backtest Performance (Net of Transaction Costs)}
\label{tab:backtest}
\resizebox{\textwidth}{!}{%
\begin{tabular}{lcccccc}
\toprule
\textbf{Strategy} & \textbf{Gross} & \textbf{Net} & \textbf{Ann.\ Ret.} & \textbf{Max DD} & \textbf{Turnover} & \textbf{Cost} \\
 & \textbf{Sharpe} & \textbf{Sharpe} & \textbf{(Net)} & \textbf{(Net)} & \textbf{(Avg)} & \textbf{(bps)} \\
\midrule
Equal Weight      & 1.215 & \textbf{1.211} & 3.9\% & $-$3.2\% & 5.0\%  & 55  \\
Risk Parity       & 1.160 & 1.155          & 3.6\% & $-$3.1\% & 5.6\%  & 62  \\
MVO + BL          & 1.032 & 1.001          & 5.2\% & $-$6.4\% & 27.0\% & 300 \\
Max Sharpe + BL   & 0.861 & 0.833          & 5.4\% & $-$4.4\% & 44.1\% & 490 \\
IC-Wt + BL        & 0.839 & 0.815          & 3.6\% & $-$7.8\% & 18.2\% & 202 \\
MVO               & 0.795 & 0.760          & 7.0\% & $-$13.3\%& 56.4\% & 626 \\
Max Sharpe        & 0.756 & 0.719          & 5.7\% & $-$7.4\% & 56.1\% & 622 \\
IC-Weighted       & 0.714 & 0.686          & 3.1\% & $-$9.0\% & 19.3\% & 214 \\
\bottomrule
\end{tabular}}
\end{table}

\textbf{Key findings:}
\begin{itemize}[noitemsep]
    \item Equal Weight and Risk Parity are the most robust---low turnover ($< 6\%$) translates to minimal transaction cost drag ($\sim$55--62 bps total over the backtest). Equal Weight's net Sharpe (1.21) barely deteriorates from gross (1.22).
    \item \textbf{BL improves MVO}: Net Sharpe rises from 0.76 to 1.00 (+0.24), while turnover drops from 56\% to 27\%. BL shrinks extreme MVO weights toward equilibrium, acting as a natural regularizer.
    \item \textbf{BL improves Max Sharpe}: Net Sharpe rises from 0.72 to 0.83 (+0.11), turnover drops from 56\% to 44\%.
    \item \textbf{BL improves IC-Weighted}: Net Sharpe rises from 0.69 to 0.82 (+0.13), with slight turnover reduction.
    \item MVO without BL has the highest turnover (56\%) and cost (626 bps), consistent with the known instability of unconstrained mean-variance portfolios.
\end{itemize}

\subsection{BL as a Regularizer}

The Black-Litterman framework effectively regularizes optimization-based strategies by blending data-driven views with an equilibrium prior. This has two practical benefits:
\begin{enumerate}[noitemsep]
    \item \textbf{Weight stability}: BL posterior returns are shrunk toward equilibrium, reducing sensitivity to estimation error in $\boldsymbol{\mu}$.
    \item \textbf{Lower turnover}: More stable return estimates produce more stable optimal weights, reducing trading costs.
\end{enumerate}

Table~\ref{tab:bl_comparison} isolates the BL effect for each optimizer:

\begin{table}[H]
\centering
\caption{Black-Litterman Impact on Optimizer Performance}
\label{tab:bl_comparison}
\begin{tabular}{lcccc}
\toprule
\textbf{Optimizer} & \textbf{Base Sharpe} & \textbf{BL Sharpe} & \textbf{$\Delta$ Sharpe} & \textbf{Turnover Reduction} \\
\midrule
MVO         & 0.760 & 1.001 & +0.241 & 56.4\% $\to$ 27.0\% \\
Max Sharpe  & 0.719 & 0.833 & +0.114 & 56.1\% $\to$ 44.1\% \\
IC-Weighted & 0.686 & 0.815 & +0.129 & 19.3\% $\to$ 18.2\% \\
\bottomrule
\end{tabular}
\end{table}

BL benefits all three optimizers, with the largest improvement for MVO (+0.24 Sharpe, turnover halved). The mechanism is return shrinkage: BL posterior returns are a precision-weighted average of the equilibrium prior and data-driven views, dampening the extreme expected returns that drive MVO instability.

\subsection{Adaptive Factor Re-Selection}

We test whether annually re-selecting factors from the full 20-factor universe (using greedy forward selection on the trailing window) improves upon the fixed factor set selected in-sample. Re-selection occurs each December; weights are rebalanced quarterly.

\begin{table}[H]
\centering
\caption{Adaptive vs Fixed Factor Selection}
\label{tab:adaptive}
\resizebox{\textwidth}{!}{%
\begin{tabular}{lcccc}
\toprule
\textbf{Strategy} & \textbf{Fixed Sharpe} & \textbf{Adaptive Sharpe} & \textbf{Factor Changes} & \textbf{Adaptive Cost} \\
\midrule
Equal Weight  & 1.211 & 0.727 & 3 & 37 bps \\
IC-Weighted   & 0.686 & 0.619 & 3 & 60 bps \\
MVO + BL      & 1.001 & 0.580 & 3 & 103 bps \\
\bottomrule
\end{tabular}}
\end{table}

Fixed factors outperform across all strategies. The adaptive approach suffers from:
\begin{enumerate}[noitemsep]
    \item \textbf{Outer turnover}: Changing factors entirely (not just reweighting) incurs large position unwinds.
    \item \textbf{Signal instability}: Factor rankings shift based on recent performance, introducing a ``factor chasing'' effect analogous to performance chasing in fund selection.
    \item \textbf{Short lookback}: Annual re-selection uses only 36 months of data, which may be insufficient to distinguish signal from noise.
\end{enumerate}

This result supports the practice of selecting factors based on long-term economic rationale rather than adaptive signal mining.

% ═══════════════════════════════════════════════════════════════
\section{Risk Analysis}
\label{sec:risk}
% ═══════════════════════════════════════════════════════════════

\subsection{Drawdown Analysis}

Our market-neutral portfolios exhibit substantially lower drawdowns than long-only benchmarks. Equal Weight's maximum drawdown is 3.3\% and IC-Weighted's is 2.2\%, compared to 14.1\% for the S\&P~500 and 21.1\% for the equal-weight mutual fund basket. This is a direct consequence of market neutrality: our returns are driven by cross-sectional dispersion rather than market direction. MAXSER variants have higher drawdowns (13--14\%) due to stock-level exposure, but remain comparable to or below the S\&P~500.

\subsection{Tail Risk}

Table~\ref{tab:oos} includes both historical VaR$_{95}$ and Cornish-Fisher VaR$_{95}$. For most of our factor portfolios, the two measures are similar, indicating approximately symmetric return distributions. Benchmark mutual funds and the S\&P~500 show larger gaps due to negative skewness and excess kurtosis---the CF adjustment captures downside risk that Gaussian VaR understates.

% ═══════════════════════════════════════════════════════════════
\section{Conclusion}
% ═══════════════════════════════════════════════════════════════

We present a complete, reproducible pipeline for cross-sectional factor investing that takes raw data as input and produces optimized, backtested portfolios as output. Our main findings:

\begin{enumerate}
    \item \textbf{Five robust factors} (AccrualRatio, CFTP, STReversal, AssetGrowth, ROE) survive multi-method selection and deliver consistent OOS alpha. All nine portfolio variants beat the S\&P~500 on risk-adjusted terms (Sharpe $\geq 0.94$ vs 0.91).

    \item \textbf{Simple allocations are hard to beat}: Equal Weight (Sharpe 1.20) outperforms all optimized factor-only variants in the static OOS test and dominates in the rolling backtest after costs (net Sharpe 1.21). This echoes \citet{demiguel2009optimal}.

    \item \textbf{MAXSER captures idiosyncratic alpha}: Factor+stock portfolios via MAXSER sparse regression (Sharpe 1.13, 14.7\% return) achieve the highest absolute returns by combining factor allocation with stock-level residual alpha, without the double-counting problem of plug-in optimizers.

    \item \textbf{Black-Litterman regularizes optimization}: BL improves MVO's net Sharpe by +0.24 and reduces turnover from 56\% to 27\%, demonstrating its value as a return shrinkage estimator. BL also improves Max Sharpe (+0.11) and IC-Weighted (+0.13).

    \item \textbf{Factor stability matters}: Fixed factor selection outperforms annual adaptive re-selection across all strategies (EW: 1.21 vs 0.73), supporting the use of economically motivated, stable factor portfolios over data-mined alternatives.

    \item \textbf{Transaction costs are material}: MVO's gross Sharpe of 0.80 falls to 0.76 net of 626 bps total cost from 56\% average turnover. Equal Weight loses only 55 bps total. Net alpha after costs is all that matters.
\end{enumerate}

\subsection{Why Equal Weight Dominates After Costs}

Equal Weight's dominance in the rolling backtest is not accidental. Quarterly rebalancing to $1/K$ systematically sells factors that appreciated (overweighted by drift) and buys those that declined, capturing a weak mean-reversion signal across factor QSpreads---the same mechanism documented by \citet{demiguel2009optimal}. Combined with only 5\% average turnover (55 bps total cost), Equal Weight preserves nearly all of its gross alpha (net Sharpe 1.21 vs gross 1.22). More sophisticated optimizers like MVO generate higher gross Sharpe ratios but surrender much of this through 56\% turnover and 626 bps in cumulative trading costs.

\subsection{Systematic Discipline and the Cost of Factor Chasing}

Our adaptive re-selection experiment directly quantifies the cost of second-guessing a systematic process. Annually re-selecting factors from the full universe based on trailing performance---which is precisely what discretionary override looks like in practice---reduces EW Sharpe from 1.21 to 0.73. This is analogous to the well-documented performance-chasing behavior in fund selection: factors that looked best over the trailing window were often mean-reverting, and the additional turnover from changing factor sets compounded the problem. The lesson for practitioners is clear: commit to the systematic framework ex ante and follow it, rather than adapting to recent performance.

\subsection{Limitations and Future Work}

Several extensions would strengthen the analysis. First, \textbf{bootstrap confidence intervals} (block bootstrap with 12-month blocks) would address whether Sharpe ratio differences are statistically meaningful given only 60 OOS months. Second, \textbf{factor risk attribution} would decompose monthly P\&L into factor contributions, answering ``where did the return come from?'' Third, \textbf{turnover-constrained optimization} would explicitly trade off Sharpe against rebalancing cost in the optimizer objective, bridging the gap between gross and net performance. Finally, \textbf{regime analysis} would reveal which factors are vulnerable to market drawdowns, addressing the practical concern of factor crowding during stress periods.

The code, data pipeline, and all results are available at:\\ \url{https://github.com/zianwang123/quant-factor-portfolio-pipeline}.

% ═══════════════════════════════════════════════════════════════
\begin{thebibliography}{99}

\bibitem[Ao et~al.(2019)]{ao2019approaching}
Ao, M., Li, Y., and Zheng, X. (2019).
\newblock Approaching mean-variance efficiency for large portfolios.
\newblock \textit{The Review of Financial Studies}, 32(7):2890--2919.

\bibitem[Carhart(1997)]{carhart1997persistence}
Carhart, M.~M. (1997).
\newblock On persistence in mutual fund performance.
\newblock \textit{The Journal of Finance}, 52(1):57--82.

\bibitem[DeMiguel et~al.(2009)]{demiguel2009optimal}
DeMiguel, V., Garlappi, L., and Uppal, R. (2009).
\newblock Optimal versus naive diversification: How inefficient is the 1/N portfolio strategy?
\newblock \textit{The Review of Financial Studies}, 22(5):1915--1953.

\bibitem[Fama and French(1993)]{fama1993common}
Fama, E.~F. and French, K.~R. (1993).
\newblock Common risk factors in the returns on stocks and bonds.
\newblock \textit{Journal of Financial Economics}, 33(1):3--56.

\bibitem[Harvey et~al.(2016)]{harvey2016and}
Harvey, C.~R., Liu, Y., and Zhu, H. (2016).
\newblock ...and the cross-section of expected returns.
\newblock \textit{The Review of Financial Studies}, 29(1):5--68.

\bibitem[He and Litterman(1999)]{he1999intuition}
He, G. and Litterman, R. (1999).
\newblock The intuition behind Black-Litterman model portfolios.
\newblock \textit{Goldman Sachs Asset Management Working Paper}.

\bibitem[Hou et~al.(2020)]{hou2020replicating}
Hou, K., Xue, C., and Zhang, L. (2020).
\newblock Replicating anomalies.
\newblock \textit{The Review of Financial Studies}, 33(5):2019--2133.

\bibitem[Ledoit and Wolf(2004)]{ledoit2004well}
Ledoit, O. and Wolf, M. (2004).
\newblock A well-conditioned estimator for large-dimensional covariance matrices.
\newblock \textit{Journal of Multivariate Analysis}, 88(2):365--411.

\end{thebibliography}

\end{document}
